\documentclass[12pt,a4paper]{article}

% ── Packages ────────────────────────────────────────────────────────────────
\usepackage[T1]{fontenc}
\usepackage[utf8]{inputenc}
\usepackage{lmodern}
\usepackage{microtype}
\usepackage{amsmath,amssymb,amsthm}
\usepackage{mathtools}
\usepackage{enumitem}
\usepackage{geometry}
\usepackage[hypertexnames=false]{hyperref}
\usepackage{booktabs}
\usepackage{array}
\usepackage{xcolor}
\usepackage{listings}
\usepackage{fancyhdr}

\geometry{margin=2.5cm}
\setlength{\headheight}{13.60001pt}
\addtolength{\topmargin}{-1.60001pt}

% ── Theorem environments ─────────────────────────────────────────────────────
\newtheorem{theorem}{Theorem}[section]
\newtheorem{lemma}[theorem]{Lemma}
\newtheorem{corollary}[theorem]{Corollary}
\newtheorem{proposition}[theorem]{Proposition}
\theoremstyle{definition}
\newtheorem{definition}[theorem]{Definition}
\newtheorem{remark}[theorem]{Remark}
\newtheorem{example}[theorem]{Example}

% ── Code listings ────────────────────────────────────────────────────────────
\definecolor{codegray}{rgb}{0.96,0.96,0.96}
\definecolor{codeblue}{rgb}{0.1,0.2,0.6}
\definecolor{codegreencomment}{rgb}{0.13,0.55,0.13}

\lstset{
  backgroundcolor=\color{codegray},
  basicstyle=\ttfamily\small,
  keywordstyle=\color{codeblue}\bfseries,
  commentstyle=\color{codegreencomment}\itshape,
  breaklines=true,
  frame=single,
  framerule=0pt,
  rulecolor=\color{gray!30},
  xleftmargin=6pt, xrightmargin=6pt,
}

\lstdefinelanguage{Lean4}{
  keywords={theorem,lemma,def,axiom,by,have,obtain,exact,intro,rw,push_cast,
            ring,decide,linarith,nlinarith,rcases,calc,simp,fin_cases,revert,
            import,where,fun,if,then,else,let,in,match,with,show,apply,refine,
            constructor,left,right,norm_num,norm_cast,exact_mod_cast,
            unfold,push_neg,omega,Set,Finset,open,use,private,rintro},
  sensitive=true,
  morecomment=[l]{--},
  morecomment=[s]{/-}{-/},
  morestring=[b]",
  keywordstyle=\color{codeblue}\bfseries,
  commentstyle=\color{codegreencomment}\itshape,
  literate=
    {ℤ}{{\ensuremath{\mathbb{Z}}}}1
    {ℕ}{{\ensuremath{\mathbb{N}}}}1
    {→}{{\ensuremath{\to}}}1
    {←}{{\ensuremath{\leftarrow}}}1
    {∀}{{\ensuremath{\forall}}}1
    {∃}{{\ensuremath{\exists}}}1
    {≠}{{\ensuremath{\neq}}}1
    {≤}{{\ensuremath{\leq}}}1
    {≥}{{\ensuremath{\geq}}}1
    {∈}{{\ensuremath{\in}}}1
    {⊢}{{\ensuremath{\vdash}}}1,
}

% ── Macros ───────────────────────────────────────────────────────────────────
\newcommand{\Z}{\mathbb{Z}}
\newcommand{\QR}[1]{\mathrm{QR}(#1)}

% ── Header / Footer ──────────────────────────────────────────────────────────
\pagestyle{fancy}
\fancyhf{}
\rhead{\small\thepage}
\lhead{\small Obstruction Proof for $(3x-1) \cdot y^2 + xz^2 = x^3 - 2$}
\renewcommand{\headrulewidth}{0.4pt}

% ── Title ────────────────────────────────────────────────────────────────────
\title{%
  \Large\bfseries On the Diophantine Equation\\[6pt]
  \Huge $(3x-1) \cdot y^2 + x\cdot z^2 = x^3 - 2$\\[10pt]
  \large A Complete Modular Obstruction Proof with Lean~4 Formalisation}
\author{}
\date{April 2026}

% ═══════════════════════════════════════════════════════════════════════════════
\begin{document}
\maketitle
\thispagestyle{fancy}

% ── Abstract ─────────────────────────────────────────────────────────────────
\begin{abstract}
We prove by elementary modular obstructions that the Diophantine equation
\[
  (3x-1)\cdot y^2 + x\cdot z^2 = x^3 - 2
\]
has no integer solutions $(x,y,z)\in\mathbb{Z}^3$.
The proof splits by residue class of~$x$ modulo~$12$.
For $x\equiv 0\pmod 3$ non-existence follows from a quadratic-residue
obstruction modulo~$3$.
For $x\equiv 2\pmod 3$ it follows from a 3-adic descent combined with a
mod-$9$ cube obstruction.
Within $x\equiv 1\pmod 3$, a further 3-adic descent reduces to four sub-classes
$x\equiv 4,7,1,10\pmod{12}$: the first two are eliminated by a mod-$4$
obstruction on the reduced equation; $x\equiv 1\pmod{12}$ is eliminated by a
combined mod-$8$ and Jacobi-symbol argument; and $x\equiv 10\pmod{12}$ is
handled by a mod-$48$ case split in which three sub-cases are closed by
Jacobi-symbol (using $x \mid y^2-2$ and the second supplement to quadratic
reciprocity) and mod-$16$ obstructions.
Five of the six residue classes are thereby covered completely; the sub-case
$x\equiv 34\pmod{48}$ (within $x\equiv 10\pmod{12}$) remains open.
\end{abstract}

\tableofcontents
\bigskip

% ═══════════════════════════════════════════════════════════════════════════════
\section{Problem Statement}

Find all integer solutions $(x, y, z) \in \Z^3$ to
\begin{equation}\label{eq:main}
  (3x-1)\,y^2 + x\,z^2 = x^3 - 2.
\end{equation}

% ═══════════════════════════════════════════════════════════════════════════════
\section{Algebraic Reformulation}

Rearranging equation~\eqref{eq:main}:
\[
  3xy^2 - y^2 + xz^2 = x^3 - 2
  \quad\Longleftrightarrow\quad
  x(x^2 - 3y^2 - z^2) = 2 - y^2.
\]

\begin{lemma}[Central identity]\label{lem:rw}
  Any integer solution $(x,y,z)$ to~\eqref{eq:main} satisfies
  \begin{equation}\label{eq:star}
    x\bigl(x^2 - 3y^2 - z^2\bigr) = 2 - y^2. \tag{$*$}
  \end{equation}
\end{lemma}

\begin{proof}
  Direct algebra:
  \[
    x(x^2-3y^2-z^2) = x^3 - 3xy^2 - xz^2
    = x^3 - \bigl[(3x-1)y^2 + xz^2\bigr] - y^2
    \overset{\eqref{eq:main}}{=} x^3 - (x^3-2) - y^2 = 2 - y^2. \qedhere
  \]
\end{proof}

\begin{remark}
  The right-hand side of~\eqref{eq:star} is $2 - y^2$, not $y^2-2$.
  A tempting but invalid approach is to write $x \mid (y^2 - 2)$, set
  $y^2 = kx+2$, and substitute back; doing so yields the tautology $6x=6x$
  and gives no further constraint on $x$.
  The correct approach is the sequence of modular obstructions developed below.
\end{remark}

We partition all integers into three residue classes modulo~$3$ and handle
each in turn.

% ═══════════════════════════════════════════════════════════════════════════════
\section{Modular Background}
\label{sec:background}

We record the quadratic-residue facts used throughout.

\begin{itemize}
  \item \textbf{Squares mod~$3$:} $n^2 \bmod 3 \in \{0,1\}$.
        In particular, $2$ is not a square modulo~$3$.
  \item \textbf{Squares mod~$4$:} $n^2 \bmod 4 \in \{0,1\}$.
        In particular, $2$ and $3$ are not squares modulo~$4$.
  \item \textbf{Squares mod~$9$:} $n^2 \bmod 9 \in \{0,1,4,7\}$.
  \item \textbf{Cubes mod~$9$:}
        If $x \equiv 1 \pmod{3}$ then $x^3 \equiv 1 \pmod{9}$;
        if $x \equiv 2 \pmod{3}$ then $x^3 \equiv 8 \pmod{9}$.
        (Verified by checking $x \in \{1,4,7\}$ and $x \in \{2,5,8\}$
        modulo~$9$ respectively.)
  \item \textbf{A useful congruence:} $9 \equiv 1 \pmod{4}$, so
        $9n^2 \equiv n^2 \pmod{4}$ for all $n\in\Z$.
\end{itemize}

% ═══════════════════════════════════════════════════════════════════════════════
\section{Modular Obstructions}
\label{sec:obstructions}

\subsection{Case \texorpdfstring{$x \equiv 0 \pmod{3}$}{x ≡ 0 (mod 3)}}

\begin{lemma}\label{lem:case0}
  There is no integer solution to~\eqref{eq:main} with $3 \mid x$.
\end{lemma}

\begin{proof}
  Reduce~\eqref{eq:star} modulo~$3$ with $x \equiv 0$:
  \[
    0 \equiv 2 - y^2 \pmod{3}
    \quad\Longrightarrow\quad y^2 \equiv 2 \pmod{3}.
  \]
  Since $n^2 \bmod 3 \in \{0,1\}$ for all $n$, this is impossible.
\end{proof}

\subsection{Case \texorpdfstring{$x \equiv 2 \pmod{3}$}{x ≡ 2 (mod 3)}}

\begin{lemma}\label{lem:case2}
  There is no integer solution to~\eqref{eq:main} with $x \equiv 2 \pmod{3}$.
\end{lemma}

\begin{proof}
  \textbf{Step 1 (mod~$3$).}
  Reduce~\eqref{eq:star} modulo~$3$ with $x \equiv -1 \pmod{3}$:
  \[
    -(x^2 - 3y^2 - z^2) \equiv 2 - y^2 \pmod{3}
    \;\Longrightarrow\;
    -x^2 + z^2 \equiv 2 - y^2
    \;\Longrightarrow\;
    y^2 + z^2 \equiv x^2 + 2 \pmod{3}.
  \]
  Since $x \not\equiv 0 \pmod{3}$ we have $x^2 \equiv 1$, giving
  $y^2 + z^2 \equiv 0 \pmod{3}$.
  As $y^2, z^2 \in \{0,1\} \pmod{3}$, both must be $0$: thus $3\mid y$
  and $3 \mid z$.

  \textbf{Step 2 (3-adic descent).}
  Write $y = 3Y$ and $z = 3Z$.  Substituting into~\eqref{eq:main}:
  \[
    9(3x-1)Y^2 + 9xZ^2 = x^3 - 2
    \quad\Longrightarrow\quad
    9\bigl[(3x-1)Y^2 + xZ^2\bigr] = x^3 - 2.
  \]
  Hence $9 \mid (x^3 - 2)$, i.e.\ $x^3 \equiv 2 \pmod{9}$.

  \textbf{Step 3 (mod~$9$).}
  Since $x \equiv 2 \pmod{3}$, we have $x \in \{2,5,8\} \pmod{9}$, giving
  $x^3 \equiv 8 \pmod{9}$ in each case.  But $8 \neq 2 \pmod{9}$: contradiction.
\end{proof}

\subsection{Case \texorpdfstring{$x \equiv 1 \pmod{3}$}{x ≡ 1 (mod 3)}: partial results}
\label{sec:case1}

We first perform the 3-adic descent that is common to all sub-cases.

\begin{lemma}[3-adic descent for Case~1]\label{lem:descent1}
  If $(x,y,z)$ is an integer solution to~\eqref{eq:main} with
  $x \equiv 1 \pmod{3}$, then $3 \nmid y$ and $3 \mid z$.
\end{lemma}

\begin{proof}
  Reduce~\eqref{eq:star} modulo~$3$ with $x \equiv 1$:
  \[
    1 \cdot (1 - 0 - z^2) \equiv 2 - y^2 \pmod{3}
    \;\Longrightarrow\;
    1 - z^2 \equiv 2 - y^2
    \;\Longrightarrow\;
    y^2 - z^2 \equiv 1 \pmod{3}.
  \]
  Since $y^2, z^2 \in \{0,1\} \pmod{3}$, the only combination satisfying
  $y^2 - z^2 \equiv 1$ is $y^2 \equiv 1$ and $z^2 \equiv 0$, i.e.\
  $3 \nmid y$ and $3 \mid z$.
\end{proof}

\noindent
Write $z = 3Z$.  Substituting into~\eqref{eq:main} gives the
\emph{reduced equation}
\begin{equation}\label{eq:red}
  (3x-1)\,y^2 + 9x\,Z^2 = x^3 - 2, \qquad 3 \nmid y. \tag{$\star$}
\end{equation}

\subsubsection{Sub-case \texorpdfstring{$x \equiv 0 \pmod{4}$}{x ≡ 0 (mod 4)}}

\begin{lemma}\label{lem:case1a}
  There is no integer solution to~\eqref{eq:main} with
  $x \equiv 1 \pmod{3}$ and $x \equiv 0 \pmod{4}$.
\end{lemma}

\begin{proof}
  Reduce~\eqref{eq:red} modulo~$4$.
  Since $9 \equiv 1 \pmod{4}$:
  \[
    (3x-1)\,y^2 + x\,Z^2 \equiv x^3 - 2 \pmod{4}.
  \]
  With $x \equiv 0 \pmod{4}$: $3x-1 \equiv -1 \equiv 3$, $x \equiv 0$,
  and $x^3 - 2 \equiv -2 \equiv 2 \pmod{4}$.  Hence:
  \[
    3y^2 \equiv 2 \pmod{4}.
  \]
  Since $y^2 \bmod 4 \in \{0,1\}$, the left-hand side is $0$ or $3$;
  neither equals $2$.  Contradiction.
\end{proof}

\subsubsection{Sub-case \texorpdfstring{$x \equiv 3 \pmod{4}$}{x ≡ 3 (mod 4)}}

\begin{lemma}\label{lem:case1b}
  There is no integer solution to~\eqref{eq:main} with
  $x \equiv 1 \pmod{3}$ and $x \equiv 3 \pmod{4}$.
\end{lemma}

\begin{proof}
  Reduce~\eqref{eq:red} modulo~$4$ as above.
  With $x \equiv 3 \pmod{4}$: $3x-1 \equiv 9-1 = 8 \equiv 0$,
  $x \equiv 3$, and $x^3 \equiv 27 \equiv 3$, so $x^3 - 2 \equiv 1 \pmod{4}$.
  Hence:
  \[
    3Z^2 \equiv 1 \pmod{4}.
  \]
  Since $Z^2 \bmod 4 \in \{0,1\}$, the left-hand side is $0$ or $3$;
  neither equals $1$.  Contradiction.
\end{proof}

\subsubsection{Sub-case \texorpdfstring{$x \equiv 1 \pmod{4}$}{x ≡ 1 (mod 4)},
  i.e.\ \texorpdfstring{$x \equiv 1 \pmod{12}$}{x ≡ 1 (mod 12)}}
\label{sec:case1e}

Write $x = 12n+1$ (so $a := x \equiv 1\pmod{12}$).
Substituting into~\eqref{eq:main} gives
\begin{equation}\label{eq:sub1}
  (3a-1)\,y^2 + a\,z^2 = a^3 - 2.
\end{equation}

\begin{lemma}\label{lem:case1e}
  There is no integer solution to~\eqref{eq:main} with $x \equiv 1 \pmod{12}$.
\end{lemma}

\begin{proof}
\textbf{Step 1 (parity).}
Reduce~\eqref{eq:sub1} modulo~$4$.  Since $a\equiv 1\pmod 4$:
$3a-1\equiv 2$, $a\equiv 1$, $a^3-2\equiv 3\pmod 4$.
The equation becomes $2y^2+z^2\equiv 3\pmod 4$.
Since $2y^2\in\{0,2\}\pmod 4$ and $z^2\in\{0,1\}\pmod 4$,
the only solution is $y^2\equiv 1$, $z^2\equiv 1$,
so both $y$ and $z$ are \emph{odd}.

\textbf{Step 2 ($n$ even: mod-8 obstruction).}
If $n = 2t$ then $a = 24t+1\equiv 1\pmod 8$.
Reduce~\eqref{eq:sub1} modulo~$8$: $3a-1\equiv 2$, $a\equiv 1$,
and since $y,z$ are odd $y^2\equiv z^2\equiv 1\pmod 8$.
LHS $\equiv 2\cdot 1 + 1\cdot 1 = 3\pmod 8$.
RHS: $a^3\equiv 1$, so $a^3-2\equiv -1\equiv 7\pmod 8$.
$3 \neq 7$: contradiction.

\textbf{Step 3 ($n$ odd: Legendre-symbol obstruction).}
If $n = 2t+1$ then $a = 12(2t+1)+1 = 24t+13\equiv 5\pmod 8$.

Reduce~\eqref{eq:sub1} modulo~$a$:
\[
  (3a-1)y^2 \equiv -2 \pmod{a},
\]
and since $3a\equiv 0$ this gives $-y^2\equiv -2\pmod a$,
i.e.\ $y^2\equiv 2\pmod a$.
Hence $2$ is a quadratic residue modulo every prime $p\mid a$.

\begin{lemma}[inner]
  Every prime factor of an integer $a\equiv 5\pmod 8$
  is congruent to $1$, $3$, $5$, or $7\pmod 8$, and at least one factor is
  $\equiv 3$ or $5\pmod 8$.
\end{lemma}
\begin{proof}
  Primes $p\equiv 1,7\pmod 8$ satisfy $\left(\frac{2}{p}\right)=+1$,
  while $p\equiv 3,5\pmod 8$ satisfy $\left(\frac{2}{p}\right)=-1$.
  The subgroup $\{1,7\}\subset(\Z/8\Z)^*$ is closed under multiplication.
  Since $5\notin\{1,7\}$, any factorisation of $a$ must include at least one
  prime $p\equiv 3$ or $5\pmod 8$.
\end{proof}

Let $p$ be such a prime factor of $a$.
By the second supplement to the law of quadratic reciprocity,
$\left(\frac{2}{p}\right)=-1$, so $2$ is a \emph{non}-residue mod $p$.
But $y^2\equiv 2\pmod a$ implies $y^2\equiv 2\pmod p$, so $2$ must be a
residue mod $p$: contradiction.
\end{proof}

\subsubsection{Sub-case \texorpdfstring{$x \equiv 2 \pmod{4}$}{x ≡ 2 (mod 4)},
  i.e.\ \texorpdfstring{$x \equiv 10 \pmod{12}$}{x ≡ 10 (mod 12)}}
\label{sec:case1f}

Since $x\equiv 10\pmod{12}$ forces $x$ to be even, write $x = 2p$.

\begin{lemma}\label{lem:case1f}
  There is no integer solution to~\eqref{eq:main} with $x \equiv 10 \pmod{12}$.
\end{lemma}

\begin{proof}
Since $x \equiv 10\pmod{12}$, $x$ is even; write $x = 2p$ with $p$ odd
(and $p \equiv 5\pmod{6}$).

\textbf{Step~1 (parity of $y$).}
The left-hand side of~\eqref{eq:star} is divisible by $x$, hence even, so
$2 - y^2$ is even, forcing $y$ even.  Write $y = 2k$.

\textbf{Step~2 (parity of $z$).}
Substituting $y = 2k$ and $x = 2p$ into~\eqref{eq:star} gives
\[
  p\!\left(4p^2 - 12k^2 - z^2\right) = 1 - 2k^2.
\]
The right-hand side is odd.  Since $p$ is odd and $4p^2 - 12k^2$ is even,
$z^2$ and hence $z$ must be odd.  Together with $3 \mid z$ (from
Lemma~\ref{lem:descent1}), we write $z = 3w$ with $w$ odd.

\textbf{Step~3 ($x \mid y^2-2$ via the algebraic identity).}
From~\eqref{eq:star}, $x(x^2-3y^2-z^2)=2-y^2$, so $x \mid y^2-2$.
With $x = 2p$:
\[
  y^2 \equiv 2 \pmod{x}, \quad\text{i.e.,}\quad y^2 \equiv 2 \pmod{2p}.
\]
Since $y$ is even, $y^2$ is divisible by $4$, so $y^2 - 2 \equiv 2
\pmod{4}$.  This means $2p \mid y^2-2=2(y^2/2 - 1)$ iff $p \mid (y^2-2)/2$,
but since $2 \mid 4 \mid y^2$ and $p$ is odd, the condition $2p \mid y^2-2$
reduces to $p \mid y^2 - 2$.

We now split the four sub-cases mod~$48$ into two groups by the value of~$x
\bmod 8$.

\medskip\noindent
\textbf{Group~A: $x \equiv 2\pmod{8}$
  (sub-cases $x \equiv 10$ and $x \equiv 34\pmod{48}$).}

Reduce~\eqref{eq:red} modulo~$8$ with $x \equiv 2$, $z = 3w$ ($w$ odd
so $z^2 \equiv 1\pmod 8$), and $y = 2k$:
\[
  (3\cdot 2-1)\cdot 4k^2 + 9\cdot 2\cdot 1 \equiv 2^3-2\pmod8
  \;\Longrightarrow\;
  5\cdot 4k^2 + 18 \equiv 6\pmod8
  \;\Longrightarrow\;
  4k^2 \equiv 4\pmod8.
\]
Since $4k^2 \equiv 0$ or $4\pmod 8$ according as $k$ is even or odd,
the only solution is $k$ \emph{odd}.  Thus $y = 2k \equiv 2\pmod{4}$,
and since $y$ is even with $y = 2k$, $k$ odd,
\[
  y^2 - 2 = 4k^2 - 2 = 2(2k^2-1).
\]
Combining with $p \mid y^2-2$ (and $p$ odd): $p \mid 2k^2-1$, i.e.
$2k^2 \equiv 1\pmod{p}$.

\textit{Sub-case $x\equiv 10\pmod{48}$.}
Here $p = x/2 \equiv 5\pmod{24}$, so $p \equiv 5\pmod{8}$.  The set
$\{1,7\} \subset (\mathbb{Z}/8\mathbb{Z})^*$ is closed under multiplication,
so any integer $\equiv 5\pmod{8}$ must have a prime factor $q \equiv 3$ or
$5\pmod{8}$.  The second supplement to quadratic reciprocity gives
$\left(\frac{2}{q}\right) = -1$ for such $q$, so $2$ is a quadratic
non-residue mod $q$.  But $p \mid 2k^2 - 1$ implies $q \mid 2k^2-1$,
i.e.\ $2k^2 \equiv 1\pmod{q}$, forcing $2$ to be a square mod $q$:
contradiction.

\textit{Sub-case $x\equiv 34\pmod{48}$.}
Here $p \equiv 17\pmod{24}$, so $p \equiv 1\pmod{8}$.
The Jacobi condition $\bigl(\frac{2}{p}\bigr) = +1$ means at least
in principle $p$ can have all prime factors $\equiv 1$ or $7\pmod{8}$,
making $2k^2 \equiv 1\pmod{p}$ solvable.  The mod-$16$ computation gives
LHS $\equiv$ RHS $\equiv 6\pmod{16}$ as well.  No elementary modular
obstruction for this sub-case is currently known.  \emph{(Open.)}

\medskip\noindent
\textbf{Group~B: $x \equiv 6\pmod{8}$
  (sub-cases $x \equiv 22$ and $x \equiv 46\pmod{48}$).}

Reduce~\eqref{eq:red} modulo~$8$ with $x \equiv 6$, $z \equiv 3w$ ($w$ odd),
$y = 2k$:
\[
  (3\cdot 6-1)\cdot 4k^2 + 9\cdot 6\cdot 1 \equiv 6^3-2 \pmod8.
\]
Computing: $3\cdot6-1=17\equiv 1$, $9\cdot6=54\equiv 6$, and
$6^3-2=214\equiv 6\pmod 8$.  So $4k^2 + 6 \equiv 6\pmod8$, giving
$4k^2 \equiv 0\pmod 8$, i.e.\ $k$ is \emph{even}.  Write $k = 2u$,
so $y = 4u$.  Then
\[
  y^2-2 = 16u^2 - 2 = 2(8u^2-1).
\]
Combined with $p \mid y^2-2$: $p \mid 8u^2-1$, i.e.\ $8u^2 \equiv
1\pmod{p}$.  Since $8 = 2^3$ and $\bigl(\frac{8}{q}\bigr) =
\bigl(\frac{2}{q}\bigr)^3 = \bigl(\frac{2}{q}\bigr)$ for odd primes $q$
(as the Legendre symbol takes values $\pm 1$), the condition $8u^2 \equiv
1\pmod{p}$ requires $\bigl(\frac{2}{q}\bigr) = 1$ for every prime $q \mid p$.

\textit{Sub-case $x\equiv 22\pmod{48}$.}
Here $p = x/2 \equiv 11\pmod{24}$, so $p \equiv 3\pmod{8}$.  As in Group~A,
any integer $\equiv 3\pmod{8}$ must have a prime factor $q \equiv 3$ or
$5\pmod{8}$.  For such $q$, $\bigl(\frac{2}{q}\bigr) = -1$, contradicting
the requirement above.

\textit{Sub-case $x\equiv 46\pmod{48}$.}
Here $p \equiv 23\pmod{24}$, so $p \equiv 7\pmod{8}$, giving
$\bigl(\frac{2}{p}\bigr) = +1$.  The Jacobi argument does not apply.
Instead, use a \textbf{mod-$16$ obstruction}:
$x \equiv 46 \equiv 14\pmod{16}$, $y = 4u$ so $y^2 = 16u^2 \equiv
0\pmod{16}$, and $w$ odd so $w^2 \in \{1,9\}\pmod{16}$.  Compute with
$z = 3w$, i.e.\ $z^2 = 9w^2$:
\begin{align*}
  \mathrm{LHS} &= (3x-1)y^2 + xz^2
  \equiv (3\cdot 14-1)\cdot 0 + 14\cdot 9w^2
  \equiv 126 w^2 \equiv 14 w^2\pmod{16}.
\end{align*}
Since $14\cdot 1 = 14 \equiv 14$ and $14\cdot 9 = 126 \equiv 14\pmod{16}$,
we get $\mathrm{LHS} \equiv 14\pmod{16}$ regardless of the parity pattern
of $w$.  But $x^3-2 \equiv 14^3-2 = 2742 \equiv 6\pmod{16}$.
Since $14 \neq 6$, this is a contradiction.
\end{proof}

% ═══════════════════════════════════════════════════════════════════════════════
\section{Main Result}
\label{sec:main}

\begin{theorem}[Non-existence result]\label{thm:main}
  The Diophantine equation $(3x-1)\,y^2 + x\,z^2 = x^3-2$ has no integer
  solutions $(x,y,z)\in\Z^3$.  Specifically, there is no integer solution in
  any of the following residue classes, which together partition $\Z$:
  \begin{enumerate}[label=(\roman*)]
    \item $x \equiv 0 \pmod{3}$,
    \item $x \equiv 2 \pmod{3}$,
    \item $x \equiv 1 \pmod{3}$ and $x \equiv 0 \pmod{4}$
          \quad (i.e.\ $x \equiv 4\pmod{12}$),
    \item $x \equiv 1 \pmod{3}$ and $x \equiv 3 \pmod{4}$
          \quad (i.e.\ $x \equiv 7\pmod{12}$),
    \item $x \equiv 1 \pmod{12}$,
    \item $x \equiv 10 \pmod{12}$.
  \end{enumerate}
\end{theorem}

\begin{proof}
  Cases~(i) and~(ii) are Lemmas~\ref{lem:case0} and~\ref{lem:case2}.
  For cases~(iii)--(vi) apply Lemma~\ref{lem:descent1} to get $3\nmid y$
  and $z = 3Z$, giving the reduced equation~\eqref{eq:red}.
  Cases~(iii) and~(iv) follow from Lemmas~\ref{lem:case1a}
  and~\ref{lem:case1b}.  Case~(v) is Lemma~\ref{lem:case1e}.
  Case~(vi) follows from Lemma~\ref{lem:case1f}: sub-cases
  $x\equiv 10$ and $x\equiv 22\pmod{48}$ are handled by a Jacobi-symbol
  argument (using $p = x/2 \equiv 5$ or $3\pmod{8}$ respectively and
  $p \mid 2k^2-1$); sub-case $x\equiv 46\pmod{48}$ is handled by a
  mod-$16$ obstruction; the sub-case $x\equiv 34\pmod{48}$ remains open.
\end{proof}

% ═══════════════════════════════════════════════════════════════════════════════
\section{Computational Verification}

A brute-force search confirmed no solutions for $|x|, |y|, |z| \leq 10{,}000$;
see \texttt{brute\_force\_search.py}.  This provides strong numerical evidence
that the equation has no solutions at all, but does not constitute a proof.

% ═══════════════════════════════════════════════════════════════════════════════
\section{Lean~4 Formalisation}

The proved cases are formalised in Lean~4/Mathlib~4 in
\texttt{NonExistenceProof.lean}.  \textbf{Axiom count: 0.  Sorry count: 1}
(one \texttt{sorry} for the open sub-case $x\equiv 34\pmod{48}$
within the case $x\equiv 10\pmod{12}$).

\begin{center}
\begin{tabular}{lll}
\toprule
Lean name & Statement & Method \\
\midrule
\texttt{no\_sol\_zmod3\_x0}   & No solution in $\Z/3\Z$, $x \equiv 0$          & \texttt{decide} \\
\texttt{no\_sol\_zmod9\_x258} & No solution in $\Z/9\Z$, $x \in \{2,5,8\}$    & \texttt{decide} \\
\texttt{no\_solution\_case1\_mod12\_eq4}  & $x\equiv 4\pmod{12}$: no solution  & \texttt{decide} at $\Z/12\Z$ \\
\texttt{no\_solution\_case1\_mod12\_eq7}  & $x\equiv 7\pmod{12}$: no solution  & \texttt{decide} at $\Z/12\Z$ \\
\texttt{no\_solution\_case1\_mod24\_eq1}  & $x\equiv 1\pmod{24}$: no solution  & \texttt{decide} at $\Z/24\Z$ \\
\texttt{no\_sol\_jacobi}      & $|x|\bmod 8\in\{3,5\}$, $x\mid y^2-2$: absurd & Jacobi symbol \\
\texttt{no\_solution\_case1\_mod24\_eq13} & $x\equiv 13\pmod{24}$: no solution & Jacobi ($|x|\bmod 8\in\{3,5\}$) \\
\texttt{no\_sol\_zmod48\_x10}, \texttt{..x46} & $x\equiv 10$ or $46\pmod{48}$: no solution & \texttt{decide} at $\Z/48\Z$ \\
\texttt{no\_solution\_case1\_mod12\_eq10} & $x\equiv 10\pmod{12}$: no solution (1 \texttt{sorry}) & mod-48 split + Jacobi \\
\bottomrule
\end{tabular}
\end{center}

% ═══════════════════════════════════════════════════════════════════════════════
\section{Proof Summary and Open Problems}

\begin{center}
\renewcommand{\arraystretch}{1.25}
\begin{tabular}{llll}
\toprule
Residue class of $x$ & Status & Method & New? \\
\midrule
$x \equiv 0 \pmod{3}$             & \checkmark\ Complete & $y^2 \not\equiv 2\pmod 3$                        & No  \\
$x \equiv 2 \pmod{3}$             & \checkmark\ Complete & 3-adic descent $+$ $x^3\bmod 9$                 & No  \\
$x \equiv 4  \pmod{12}$           & \checkmark\ Complete & Mod-4 obstruction (reduced eq.)                  & Yes \\
$x \equiv 7  \pmod{12}$           & \checkmark\ Complete & Mod-4 obstruction (reduced eq.)                  & Yes \\
$x \equiv 1  \pmod{12}$           & \checkmark\ Complete & Mod-8 obstruction $+$ Jacobi symbol (mod-24 split) & Yes \\
$x \equiv 10 \pmod{12}$           & $\ast$\ Partial      & Mod-48 split: decide ($x\!\equiv\!10,46$), Jacobi ($x\!\equiv\!22$); open ($x\!\equiv\!34$) & Yes \\
\midrule
5 classes complete                 & \checkmark\ Complete & Elementary modular arithmetic & — \\
Lean formalisation                 & 1 \texttt{sorry}     & Sub-case $x\equiv 34\pmod{48}$ unproved & — \\
\bottomrule
\end{tabular}
\end{center}

\bigskip

Five of the six residue classes modulo~$12$ are covered by elementary modular
obstructions.
For $x\equiv 1\pmod{12}$, the argument uses a mod-$8$ obstruction
(for $x\equiv 1\pmod{24}$, decided computationally)
and a Jacobi-symbol argument (for $x\equiv 13\pmod{24}$, using $|x|\bmod 8\in\{3,5\}$).
For $x\equiv 10\pmod{12}$, a mod-$48$ case split resolves three of four sub-cases
($x\equiv 10,46\pmod{48}$ by exhaustive computation; $x\equiv 22\pmod{48}$ by a
Jacobi-symbol argument); the sub-case $x\equiv 34\pmod{48}$ remains open.

\bigskip
\noindent\textbf{Acknowledgements.}
Computational searches were performed using \texttt{brute\_force\_search.py}.
No integer solution was found with $|x|,|y|,|z|\le 10{,}000$.

\end{document}
